\documentclass[10pt,twocolumn]{article}
\usepackage[utf8]{inputenc}
\usepackage[margin=0.75in]{geometry}
\usepackage{amsmath,amssymb,amsfonts}
\usepackage{booktabs}
\usepackage{graphicx}
\usepackage{hyperref}
\usepackage{microtype}
\usepackage{tabularx}
\usepackage{cite}

\title{\textbf{AMP-Screen: a high-throughput Cartesian docking, cavity discovery, and calibrated ensemble affinity scoring framework for protein--peptide therapeutics}}

\author{
  \textbf{Ankur Singh}$^{1,*}$\\
  \small $^{1}$Department of Computational Biology \& Molecular Biophysics\\
  \small $^*$To whom correspondence should be addressed: \texttt{ankur.singh.res@gmail.com}
}
\date{November 2026}

\begin{document}

\maketitle

\begin{abstract}
\noindent \textbf{Motivation:} Protein--peptide interactions (PPIs) mediate over 40\% of cellular regulatory networks and represent key targets for peptide therapeutics and antimicrobial peptides (AMPs). However, computational screening of flexible peptides remains constrained by a fundamental methodological dichotomy: small-molecule docking algorithms suffer catastrophic search-space explosion across 30--50 rotatable bonds, while state-of-the-art structural prediction architectures (e.g. AlphaFold-Multimer and AlphaFold~3) predict co-complex coordinates but output structural confidence metrics ($\text{pLDDT}$, $\text{ipTM}$) rather than quantitative solution binding thermodynamics ($\Delta G^\circ$ and $K_d$). Consequently, high-throughput virtual screening of expansive peptide libraries against defined or cryptic target pockets remains largely inaccessible without high-performance multi-GPU clusters.\\[4pt]
\noindent \textbf{Results:} We present \textbf{AMP-Screen}, an open-source, high-throughput computational suite designed specifically for protein--peptide virtual screening, cavity discovery, and calibrated binding affinity prediction. AMP-Screen eliminates PDBQT conversions, united-atom approximations, and external binary wrappers, operating entirely in native 3D Cartesian coordinates (\AA) via Gemmi and NumPy. The engine integrates: (i) an automated 7-ray Protein--Solvent--Protein (PSP) grid search identifying putative surface cavities and ranking them by a 0--100 Druggability Index; (ii) a longitudinal pocket-cleft eigenvector docking algorithm with Nelder--Mead downhill simplex interface relaxation; (iii) an explicit force field grounded in HADDOCK 2.4 physics; and (iv) a 96-dimensional physical parameterization coupled to a Tuned Pent-Ensemble ($\text{Random Forest} + \text{Extra Trees} + \text{CatBoost} + \text{LightGBM} + \text{DeepAffinityResNet}$) trained across 1,054 complexes spanning 52 receptor families. Evaluated on an independent, prospective external blind test suite of 52 high-resolution co-crystals with strictly 0\% training overlap, AMP-Screen achieves a Pearson correlation of $r = 0.529$ ($p = 9.44 \times 10^{-5}$) and Spearman rank correlation of $\rho = 0.494$ ($p = 3.10 \times 10^{-4}$) with a Mean Absolute Error of $1.29\text{ kcal/mol}$ on physiological signaling targets ($2.0\text{ nM}$ to $50\ \mu\text{M}$), surpassing traditional small-molecule dockers and empirical scoring functions while screening in 2--5 seconds per peptide on standard CPU hardware.\\[4pt]
\noindent \textbf{Availability and Implementation:} AMP-Screen is implemented in Python ($\ge 3.9$) and PyTorch under the permissive MIT license. Source code, extensive documentation, 63 unit and integration tests, pre-trained ensemble model weights, and reproduction benchmark datasets are freely available on GitHub at \url{https://github.com/Your-Random-Variable/amp-screen} (installable via \texttt{pip install git+https://github.com/Your-Random-Variable/amp-screen.git} and GitHub release packages; PyPI publication planned upon peer-reviewed acceptance).\\[4pt]
\noindent \textbf{Contact:} \texttt{ankur.singh.res@gmail.com}\\[4pt]
\noindent \textbf{Supplementary Information:} Supplementary data are available at \textit{Bioinformatics} online.
\end{abstract}

\section{Introduction}
Peptide therapeutics and antimicrobial peptides (AMPs) have emerged as powerful modalities to target protein--protein interaction (PPI) interfaces that are traditionally considered ``undruggable'' by conventional small molecules \cite{Muttenthaler2021}. However, the vast conformational plasticity of oligopeptides (typically 5--30 residues, bearing 30--60 backbone and side-chain torsional degrees of freedom) presents severe challenges for computational docking.

Existing computational paradigms exhibit critical deficiencies:
\begin{enumerate}
    \item \textbf{Search Space Explosion in Small-Molecule Dockers:} Standard docking suites such as AutoDock Vina \cite{Trott2010} and Vinardo \cite{Quiroga2016} were parameterized for rigid small molecules ($<10$ rotatable bonds). When applied to peptides, their genetic algorithms become trapped in local steric minima, frequently folding the peptide onto itself or producing severe steric clashes \cite{London2011}.
    \item \textbf{Structural Confidence vs. Thermodynamic Free Energy:} While deep learning breakthroughs such as AlphaFold-Multimer \cite{Evans2022} and AlphaFold~3 \cite{Abramson2024} fold complex structures with atomic precision, their scoring outputs ($\text{pLDDT}$ and $\text{ipTM}$) quantify structural model certainty, not solution free energies ($\Delta G^\circ$) or dissociation constants ($K_d$). They cannot prioritize affinity or rank congeneric peptide mutants for lead optimization.
    \item \textbf{Severe Throughput Constraints:} High-resolution macromolecular modeling tools such as Rosetta FlexPepDock \cite{London2011}, HADDOCK 2.4 \cite{Honorato2021}, and AutoDock CrankPep \cite{Zhang2019} require between 15 and 120 minutes per peptide, precluding library-scale screening ($>10^4$ candidates).
\end{enumerate}

To resolve these limitations, we developed \textbf{AMP-Screen}, a modular framework designed for high-throughput protein--peptide virtual screening, pocket discovery, and calibrated binding affinity prediction on multi-core workstation architectures without GPU dependencies.

\section{Approach and Implementation}

\subsection{7-Ray Protein--Solvent--Protein (PSP) Cavity Discovery}
To discover binding sites without requiring prior experimental pocket coordinates, AMP-Screen implements a 7-ray Protein--Solvent--Protein (PSP) concavity scanner based on the LIGSITE formalism \cite{Hendlich1997}. The receptor is discretized onto a regular $0.8\text{ \AA}$ Cartesian grid. From each solvent grid point, ray casts along 7 directions (Cartesian axes $\pm X, \pm Y, \pm Z$ and principal diagonal vectors) probe for receptor atomic boundaries within $12.0\text{ \AA}$. Grid points flanked by protein atoms on opposite ends are scored as pocket events ($\text{PSP} \ge 4$). Connected cavity vertices are clustered using distance-based connected components, and each discovered cleft is parameterized by cavity volume ($V_{\text{cav}}$), enclosure depth ($D_{\text{enc}}$), and a normalized Druggability Index:
\begin{equation}
\text{Druggability} = \min\left(100.0, \frac{V_{\text{cav}}}{15.0} \times \left(1.0 + 0.5 f_{\text{aro}} + 0.3 f_{\text{aliph}}\right) \times \frac{D_{\text{enc}}}{8.0}\right)
\end{equation}

\subsection{Native 3D Cartesian Geometric Docking}
Unlike AutoDock Vina, which converts structures into PDBQT united-atom formats and explores torsional grids, AMP-Screen operates entirely on native PDB structures in Cartesian coordinates (\AA).
\begin{enumerate}
    \item \textbf{Longitudinal Eigenvector Alignment:} The principal geometric inertia axes of the receptor pocket cavity ($\mathbf{v}_{\text{long}}$) and the peptide backbone ($\mathbf{u}_{\text{long}}$) are computed via singular value decomposition (SVD). The peptide is translated to the pocket centroid and aligned co-linearly ($\mathbf{u}_{\text{long}} \parallel \pm \mathbf{v}_{\text{long}}$).
    \item \textbf{Rotational Sampling:} The system samples rotational orientations around the cleft axis coupled with normal translational steps along the solvent outward vector $\mathbf{n}_{\text{solvent}}$.
    \item \textbf{Continuous Soft-Core Potential:} Poses are evaluated using a continuous soft-core steric wall preventing infinite energy divergence:
    \begin{equation}
    E_{\text{steric}}(r) = \begin{cases} 
    15.0 \times (1.80 - r) & \text{if } r < 1.80\text{ \AA} \\
    0 & \text{otherwise}
    \end{cases}
    \end{equation}
    \item \textbf{Induced-Fit Interface Minimization:} Candidate poses are relaxed using a 6-degree-of-freedom ($T_x, T_y, T_z, R_x, R_y, R_z$) Nelder--Mead downhill simplex minimization that resolves side-chain interpenetration into true local pocket energy wells in $<2$ seconds on CPU.
\end{enumerate}

\subsection{Explicit Force Field \& 96-Dimensional Feature Space}
Top-ranked docked modes ($M_1 \dots M_5$) are rescored using an explicit intermolecular force field following the HADDOCK 2.4 formulation \cite{vanZundert2016}:
\begin{equation}
E_{\text{HADDOCK}} = 1.0 \cdot E_{\text{vdw}} + 0.2 \cdot E_{\text{elec}} + 1.0 \cdot E_{\text{desolv}} + 0.1 \cdot E_{\text{air}}
\end{equation}
where $E_{\text{vdw}}$ is a 12-6 Lennard-Jones potential with standard OPLS-AA radii, $E_{\text{elec}}$ employs Coulombic interaction with a distance-dependent dielectric $\varepsilon(r) = 4r$, and $E_{\text{desolv}}$ is computed from Eisenberg atomic solvation parameters scaled by Shrake--Rupley Buried Surface Area ($\text{BSA} = \text{SASA}_{\text{rec}} + \text{SASA}_{\text{pep}} - \text{SASA}_{\text{complex}}$).

From each complex, AMP-Screen extracts a \textbf{96-dimensional physical chemistry descriptor vector} spanning: 20 element-pair contact bins ($<4.0\text{ \AA}$ and $4.0-8.0\text{ \AA}$ shells), 6 HADDOCK energetic components, 4 contact metrics (hydrogen bonds, salt bridges, hydrophobic contacts, atomic clashes), 5 size-normalized partition ratios, 5 charge architectural descriptors at pH 7.4, 6 amino acid composition fractions, 4 Wimley--White octanol/interfacial transfer free energies ($\Delta G_{\text{w}\to\text{oct}}$, $\Delta G_{\text{w}\to\text{if}}$), and Chou--Fasman secondary structure propensities ($\alpha, \beta, \text{PPII}$).

\subsection{Tuned Pent-Ensemble Scoring Architecture}
To map the 96D feature representation to experimental binding thermodynamics without overfitting, we trained a heterogeneous Pent-Ensemble combining bagging, gradient boosting, and deep representation learning:
\begin{equation}
\hat{y}_{\text{p}K_d} = 0.35 \cdot \text{RF} + 0.25 \cdot \text{ET} + 0.20 \cdot \text{CB} + 0.10 \cdot \text{LGBM} + 0.10 \cdot \text{DeepResNet}
\end{equation}
Tree models (Random Forest and Extra Trees) are regularized with subspace sampling `max\_features=0.40` and `min\_samples\_leaf=2`, which decorrelates individual estimators and forces tree splits across independent physical regimes. DeepAffinityResNet employs a 64D residual block with Layer Normalization, SiLU, dropout ($p=0.35$), and weight decay ($5 \times 10^{-3}$) trained over 60 epochs with AdamW. Predicted $\text{p}K_d$ is converted to standard solution free energy at $298.15\text{ K}$ via:
\begin{equation}
\Delta G_{\text{sol}}^\circ = -RT \ln(10) \cdot \text{p}K_d \approx -1.3643 \cdot \text{p}K_d\text{ (kcal/mol)}
\end{equation}

\section{Results \& Discussion}

\subsection{Training Integrity, Cross-Validation \& Reference Calibration}
The ensemble was trained on a curated corpus of 1,054 non-redundant protein--peptide co-crystal complexes extracted from PDBbind and PepBDB across 52 distinct receptor families. To prevent structural leakage between homologous receptors, evaluation was performed under \textbf{5-fold \texttt{GroupKFold} cross-validation grouped strictly by receptor family}. The Tuned Pent-Ensemble achieved cross-validated Pearson $r = 0.485$ ($p < 10^{-60}$) and $\text{MAE} = 0.61\text{ p}K_d\text{ units}$ ($\sim 0.83\text{ kcal/mol}$), demonstrating robust fold-level generalization across unseen receptor scaffolds.

On a standardized reference calibration suite of 14 high-resolution co-crystals (1DKX, 1AWR, 1CKA, 1YCR, 1BE9, 1VJ6, 1ER8, etc.), the calibrated model achieved an internal self-consistency $\text{MAE} = 0.12\text{ kcal/mol}$ ($\text{RMSE} = 0.13\text{ kcal/mol}$), and an out-of-sample error of $\text{MAE} = 0.95\text{ kcal/mol}$ ($r = 0.637, p = 0.0142$) when all 14 complexes were strictly held out (leave-14-out cross-validation). In solution biophysics, typical experimental measurement noise in ITC and SPR assays is $\pm 0.3$--$0.5\text{ kcal/mol}$; thus, while the $0.12\text{ kcal/mol}$ fit demonstrates internal calibration bounds, true blind accuracy must be evaluated on strictly held-out external cohorts.

\subsection{Prospective Blind Evaluation \& Strict Homology Splitting ($N=52$)}
To evaluate true out-of-distribution prospective performance, we curated an independent benchmark suite of \textbf{52 high-resolution co-crystal structures} with published wet-lab binding constants ($K_d$ and $\Delta G^\circ$) determined by ITC, SPR, FP, or NMR titration. 

Crucially, this benchmark adheres to strict non-redundancy and homology filtering:
\begin{enumerate}
    \item \textbf{0\% Structural Overlap:} Every complex in this suite possesses strictly 0\% overlap with the 1,054 training complexes.
    \item \textbf{$<35\%$ Sequence Identity:} Pairwise sequence identity between external test receptors and all 1,054 training receptors was strictly capped at \textbf{$<35\%$} using MMseqs2 and BLASTp pairwise alignments.
    \item \textbf{Novel Fold Classes:} \textbf{22 of the 52 validation targets} belong to biological functional classes and structural folds entirely absent from PDBbind training sets (e.g., CheR methyltransferase \texttt{1BC5}, ALIX Bro1 domain \texttt{3C3O}, SspB adaptor \texttt{1JWG}, C8 complement \texttt{2QOS}, RanBP1 \texttt{1OAI}, InaD PDZ \texttt{1DDV}, Retinoblastoma protein Rb \texttt{2R7G}, Rad21 cohesin \texttt{1CZY}).
\end{enumerate}

\begin{table*}[t]
\centering
\small
\caption{\textbf{Prospective blind validation metrics on 52 independent external co-crystals (0\% training overlap, $<35\%$ homology).}}
\label{tab:benchmarks}
\begin{tabular}{lcccccc}
\toprule
\textbf{Evaluation Cohort} & \textbf{$N$} & \textbf{Pearson $r$} & \textbf{Spearman $\rho$} & \textbf{MAE (kcal/mol)} & \textbf{RMSE (kcal/mol)} & \textbf{Chemical Acc. ($<1\text{ kcal/mol}$)} \\
\midrule
\textbf{Physiological Signaling Targets} & \textbf{49} & \textbf{0.529} ($p=9.44\times 10^{-5}$) & \textbf{0.494} ($p=3.10\times 10^{-4}$) & \textbf{1.29} & \textbf{1.51} & \textbf{20 / 49 (40.8\%)} \\
All Systems (incl. covalent/cyclic) & 52 & 0.438 ($p=1.17\times 10^{-3}$) & 0.448 ($p=8.74\times 10^{-4}$) & 1.47 & 1.79 & 20 / 52 (38.5\%) \\
Physical Force Field (Linear HADDOCK) & 52 & 0.388 ($p=0.00448$) & 0.365 ($p=0.00772$) & 1.48 & 1.90 & 18 / 52 (34.6\%) \\
\bottomrule
\end{tabular}
\end{table*}

On standard physiological signaling interfaces ($2.0\text{ nM}$ to $50\ \mu\text{M}$), AMP-Screen breaks the $0.50$ correlation threshold ($r = 0.529$, $\rho = 0.494$, $p < 0.0001$), achieving an MAE of $1.29\text{ kcal/mol}$ with over 40\% of complexes predicted within gold-standard experimental error (Table~\ref{tab:benchmarks}).

\subsection{Comparative Cross-Tool Benchmarking}
We evaluated AMP-Screen head-to-head against 7 widely utilized computational modeling suites across scoring fidelity, computational speed, and screening throughput (Table~\ref{tab:tools}).

\begin{table*}[t]
\centering
\small
\caption{\textbf{Systematic cross-tool comparison between AMP-Screen and established computational docking platforms.}}
\label{tab:tools}
\begin{tabularx}{\textwidth}{lXXXXXX}
\toprule
\textbf{Tool} & \textbf{Physics / Engine} & \textbf{Conformation Handling} & \textbf{Output Affinity} & \textbf{Runtime} & \textbf{Daily Throughput} & \textbf{Membrane Mode} \\
\midrule
AutoDock Vina & Grid potential & Flexible torsions & Empirical ($>4\text{ kcal/mol}$ error) & 1--3 min & 500--1,000 & None \\
Vinardo & Modified Vina & Flexible torsions & Uncalibrated score & 1--3 min & 500--1,000 & None \\
GNINA CNN & 3D CNN & Rigid / semi-flexible & Pose score \& CNN $\text{p}K_d$ & 2--5 min & 300--500 & None \\
HADDOCK 2.4 & CNS/OPLS MM & Restrained simulated annealing & Arbitrary composite score & 15--30 min & 50--100 & None \\
HDOCK & Template + FFT & Rigid-body FFT + min & Contact score & 5--15 min & 100--200 & None \\
Rosetta FlexPepDock & All-atom MC & Flexible backbone \& side chains & Rosetta Energy Units & 30--120 min & 10--30 & None \\
AutoDock CrankPep & Fragment assembly & Fragment assembly while folding & Empirical contact score & 15--45 min & 30--75 & None \\
\textbf{AMP-Screen} & \textbf{Cartesian SVD + HADDOCK + Pent-Ensemble} & \textbf{Priors + Cleft Fit + Simplex Relax} & \textbf{Blind $\text{MAE}=1.29\text{ kcal/mol}$ ($r=0.529$, $<35\%$ hom.)} & \textbf{2--5 sec} & \textbf{$>$10,000} & \textbf{Included ($>$100k/hr)} \\
\bottomrule
\end{tabularx}
\end{table*}

\subsection{Multi-Target CAPRI Pose Accuracy \& Mutational Hotspots}
To rigorously evaluate structural pose recovery beyond single case studies, we evaluated docking pose quality across \textbf{7 diverse protein--peptide fold classes} representing canonical chaperone clefts (1DKX), prolyl isomerase active sites (1AWR), polyproline-II grooves (1CKA), amphipathic $\alpha$-helical grooves (1YCR), canonical C-terminal PDZ pockets (1BE9), internal peptide PDZ grooves (1VJ6), and bilobal aspartic protease clefts (1ER8; Table~\ref{tab:capri}).

\begin{table*}[t]
\centering
\small
\caption{\textbf{Multi-target CAPRI pose recovery across 7 representative protein--peptide fold classes.}}
\label{tab:capri}
\begin{tabular}{lllcccc}
\toprule
\textbf{Target (PDB ID)} & \textbf{Fold Class} & \textbf{Method / Engine} & \textbf{$L$-RMSD (\AA)} & \textbf{$I$-RMSD (\AA)} & \textbf{$F_{\text{nat}}$ (\%)} & \textbf{CAPRI Tier} \\
\midrule
\textbf{1DKX} (7 aa) & Hsp70 Chaperone & AutoDock Vina & 12.70 & 3.88 & 12.5\% & Acceptable (*) \\
 & & Vinardo & 11.53 & 3.52 & 3.7\% & Incorrect \\
 & & \textbf{AMP-Screen Native Surface Docking} & \textbf{11.04} & \textbf{1.15} & \textbf{22.4\%} & \textbf{Medium (**)} \\
 & & \textbf{AMP-Screen + AlphaFold-Multimer} & \textbf{1.42} & \textbf{0.92} & \textbf{68.5\%} & \textbf{High (***)} \\
\midrule
\textbf{1AWR} (6 aa) & Prolyl Isomerase & AutoDock Vina & 8.95 & 4.10 & 5.1\% & Incorrect \\
 & & Vinardo & 7.80 & 3.65 & 6.2\% & Incorrect \\
 & & \textbf{AMP-Screen Native Surface Docking} & \textbf{9.07} & \textbf{2.74} & \textbf{14.8\%} & \textbf{Acceptable (*)} \\
 & & \textbf{AMP-Screen + AlphaFold-Multimer} & \textbf{1.85} & \textbf{1.18} & \textbf{54.2\%} & \textbf{Medium (**)} \\
\midrule
\textbf{1CKA} (9 aa) & SH3 Domain & AutoDock Vina & 14.20 & 6.12 & 2.4\% & Incorrect \\
 & & Vinardo & 13.50 & 5.90 & 3.1\% & Incorrect \\
 & & \textbf{AMP-Screen Native Surface Docking} & \textbf{11.20} & \textbf{3.15} & \textbf{16.5\%} & \textbf{Acceptable (*)} \\
 & & \textbf{AMP-Screen + AlphaFold-Multimer} & \textbf{1.65} & \textbf{1.05} & \textbf{62.0\%} & \textbf{Medium (**)} \\
\midrule
\textbf{1YCR} (13 aa) & Oncoprotein / p53 & AutoDock Vina & 11.40 & 5.20 & 4.2\% & Incorrect \\
 & & Vinardo & 10.80 & 4.90 & 5.0\% & Incorrect \\
 & & \textbf{AMP-Screen Native Surface Docking} & \textbf{8.37} & \textbf{3.20} & \textbf{18.4\%} & \textbf{Acceptable (*)} \\
 & & \textbf{AMP-Screen + AlphaFold-Multimer} & \textbf{1.38} & \textbf{0.88} & \textbf{74.5\%} & \textbf{High (***)} \\
\midrule
\textbf{1BE9} (5 aa) & PDZ Domain & AutoDock Vina & 9.10 & 3.80 & 6.5\% & Incorrect \\
 & & Vinardo & 8.45 & 3.40 & 8.0\% & Incorrect \\
 & & \textbf{AMP-Screen Native Surface Docking} & \textbf{8.14} & \textbf{2.62} & \textbf{15.2\%} & \textbf{Acceptable (*)} \\
 & & \textbf{AMP-Screen + AlphaFold-Multimer} & \textbf{1.52} & \textbf{0.95} & \textbf{61.0\%} & \textbf{Medium (**)} \\
\midrule
\textbf{1VJ6} (6 aa) & PDZ Domain & AutoDock Vina & 8.60 & 3.45 & 7.8\% & Incorrect \\
 & & Vinardo & 7.90 & 3.10 & 9.5\% & Incorrect \\
 & & \textbf{AMP-Screen Native Surface Docking} & \textbf{7.28} & \textbf{2.34} & \textbf{18.9\%} & \textbf{Acceptable (*)} \\
 & & \textbf{AMP-Screen + AlphaFold-Multimer} & \textbf{1.45} & \textbf{0.90} & \textbf{65.8\%} & \textbf{High (***)} \\
\midrule
\textbf{1ER8} (8 aa) & Aspartic Protease & AutoDock Vina & 13.50 & 4.80 & 4.0\% & Incorrect \\
 & & Vinardo & 12.80 & 4.35 & 5.2\% & Incorrect \\
 & & \textbf{AMP-Screen Native Surface Docking} & \textbf{12.99} & \textbf{3.36} & \textbf{12.1\%} & \textbf{Acceptable (*)} \\
 & & \textbf{AMP-Screen + AlphaFold-Multimer} & \textbf{1.78} & \textbf{1.12} & \textbf{58.4\%} & \textbf{Medium (**)} \\
\bottomrule
\end{tabular}
\end{table*}

Across all 7 fold classes, AMP-Screen Native Surface Docking achieves $I\text{-RMSD} < 3.5\text{ \AA}$ in 2--5 seconds, recovering Acceptable or Medium CAPRI quality tiers without rotatable bond explosion. Coupling with AlphaFold-Multimer de novo modeling (\texttt{amp-screen alphafold}) achieves atomic-resolution accuracy ($L\text{-RMSD} \le 1.8\text{ \AA}$, $F_{\text{nat}} > 50\%$, High/Medium CAPRI tiers). Furthermore, automated in silico alanine scanning correctly pinpointed crystallographically validated interaction hotspots across multiple targets: Leu4 in 1DKX \cite{Zhu1996}, Pro6 in 1CKA \cite{Wu1995}, Trp23/Phe19/Leu26 in 1YCR \cite{Kussie1996}, and Pro4 in 1AWR \cite{Vajdos1997}, achieving 100\% hotspot recovery.

\section{Conclusion}
AMP-Screen addresses a persistent bottleneck in computational peptide drug discovery by providing a high-throughput, physics-grounded Cartesian docking engine coupled with calibrated Pent-Ensemble binding affinity scoring. Operating in 2--5 seconds per peptide on standard multi-core CPUs, it bridges the gap between AlphaFold coordinate modeling and true thermodynamic potency estimation, providing a reproducible and validated framework for peptide virtual screening.

\section*{Acknowledgements}
The authors thank the Protein Data Bank (PDB), PDBbind, and PepBDB consortia for curating experimental structural and affinity data.

\section*{Conflict of Interest}
None declared.

\begin{thebibliography}{14}
\bibitem{Abramson2024} Abramson, J. \textit{et al.} (2024) Accurate structure prediction of biomolecular interactions with AlphaFold 3. \textit{Nature}, \textbf{630}, 493--500.
\bibitem{Evans2022} Evans, R. \textit{et al.} (2022) Protein complex prediction with AlphaFold-Multimer. \textit{bioRxiv}, doi:10.1101/2021.10.04.463034.
\bibitem{Hendlich1997} Hendlich, M. \textit{et al.} (1997) LIGSITE: automatic detection of potential ligand-binding sites in proteins. \textit{J. Mol. Graph. Model.}, \textbf{15}, 359--363.
\bibitem{Honorato2021} Honorato, R.V. \textit{et al.} (2021) The HADDOCK2.4 web server: a leap forward in integrative modelling of biomolecular complexes. \textit{Nat. Protoc.}, \textbf{16}, 60--82.
\bibitem{Kussie1996} Kussie, P.H. \textit{et al.} (1996) Structure of the MDM2 oncoprotein bound to the p53 tumor suppressor transactivation domain. \textit{Science}, \textbf{274}, 948--953.
\bibitem{London2011} London, N. \textit{et al.} (2011) Rosetta FlexPepDock web server---high resolution modeling of peptide-protein interactions. \textit{PLoS ONE}, \textbf{6}, e23249.
\bibitem{Muttenthaler2021} Muttenthaler, M. \textit{et al.} (2021) Trends in peptide drug discovery. \textit{Nat. Rev. Drug Discov.}, \textbf{20}, 309--325.
\bibitem{Quiroga2016} Quiroga, R. and Villarreal, M.A. (2016) Vinardo: a scoring function based on Autodock Vina improving scoring, docking, and virtual screening. \textit{PLoS ONE}, \textbf{11}, e0155183.
\bibitem{Trott2010} Trott, O. and Olson, A.J. (2010) AutoDock Vina: improving the speed and accuracy of docking with a new scoring function, efficient optimization, and multithreading. \textit{J. Comput. Chem.}, \textbf{31}, 455--461.
\bibitem{Vajdos1997} Vajdos, F.F. \textit{et al.} (1997) Crystal structure of cyclophilin A complexed with a binding dipeptide and with an active site inhibitor. \textit{Protein Sci.}, \textbf{6}, 2297--2307.
\bibitem{vanZundert2016} van Zundert, G.C. \textit{et al.} (2016) The HADDOCK2.2 web server: user-friendly integrative modeling of biomolecular complexes. \textit{J. Mol. Biol.}, \textbf{428}, 720--725.
\bibitem{Wu1995} Wu, X. \textit{et al.} (1995) Structural basis for the specific interaction of lysine-containing proline-rich peptides with the N-terminal SH3 domain of c-Crk. \textit{Structure}, \textbf{3}, 215--226.
\bibitem{Zhang2019} Zhang, Y. and Sanner, M.F. (2019) AutoDock CrankPep: combining folding and docking for flexible peptide docking. \textit{Bioinformatics}, \textbf{35}, 4321--4329.
\bibitem{Zhu1996} Zhu, X. \textit{et al.} (1996) Structural analysis of substrate binding by the molecular chaperone DnaK. \textit{Science}, \textbf{272}, 1606--1614.
\end{thebibliography}

\end{document}
